\documentclass[a4paper,11pt]{article}
\usepackage[T2A]{fontenc}
\usepackage[utf8]{inputenc}
\usepackage[russian]{babel}
\usepackage{amsmath,amssymb,mathtools}
\usepackage{DejaVuSans}
\renewcommand{\familydefault}{\sfdefault}
\usepackage{geometry}
\geometry{left=24mm,right=24mm,top=23mm,bottom=23mm}
\usepackage{microtype,icomma}
\usepackage{booktabs,tabularx}
\usepackage{enumitem}
\usepackage{xcolor,tikz}
\usetikzlibrary{arrows.meta,positioning,shapes.geometric,calc}
\usepackage{fancyhdr,setspace,hyperref}
\hypersetup{colorlinks=true,linkcolor=teal!65!black,urlcolor=teal!65!black}
\definecolor{accent}{HTML}{087E82}
\definecolor{darkaccent}{HTML}{075B60}
\definecolor{ink}{HTML}{292D31}
\definecolor{muted}{HTML}{617075}
\color{ink}
\setstretch{1.12}
\setlength{\parindent}{0pt}
\setlength{\parskip}{3pt}
\setlist[enumerate]{leftmargin=8mm,itemsep=4pt,topsep=3pt}
\setlist[itemize]{leftmargin=7mm,itemsep=2pt,topsep=3pt}
\setlength{\headheight}{15pt}
\pagestyle{fancy}
\fancyhf{}
\fancyhead[L]{\small\color{muted}Проценты, доли и нахождение целого}
\fancyhead[R]{\small\color{muted}Надя Клименко}
\fancyfoot[C]{\small\color{muted}\thepage}
\renewcommand{\headrulewidth}{0.35pt}
\renewcommand{\footrulewidth}{0pt}
\newcommand{\mysection}[2]{\section*{\color{darkaccent}#1. #2}\vspace{-3pt}{\color{accent}\hrule height 0.55pt}\vspace{7pt}}
\newcommand{\idea}[1]{\par\smallskip{\color{darkaccent}\bfseries #1}\par\smallskip}
\newcommand{\ex}[1]{\par\medskip{\color{darkaccent}\bfseries Пример.} #1\par\smallskip}
\newcommand{\note}[1]{\par\smallskip{\color{accent}\bfseries Важно.} #1\par\smallskip}
\begin{document}
\hypersetup{pageanchor=false}
\begin{titlepage}
\thispagestyle{empty}
\begin{center}
\vspace*{27mm}
{\large\color{darkaccent}\textbf{МАТЕМАТИКА}}\par
\vspace{22mm}
{\LARGE Чек-лист}\par
\vspace{12mm}
{\LARGE\color{accent}\bfseries Проценты, доли\\[3mm]и нахождение целого}\par
\vspace{17mm}
{\large 6 класс \quad\textbullet\quad Математическая вертикаль}\par
\vfill
{\large Надя Клименко}\par
\vspace{12mm}
{\normalsize 7 ноября 2026 г.}\par
\vspace{22mm}
{\small Лёвин Артём Александрович\\[1mm]эксклюзивно для Нади Клименко}\par
\vspace{15mm}
\begin{tikzpicture}[x=1cm,y=1cm]
\draw[accent,line width=0.8pt] (0,0) .. controls (1.4,0.6) and (2.4,-0.5) .. (3.6,0)
 .. controls (4.8,0.6) and (5.8,-0.5) .. (7.2,0);
\end{tikzpicture}
\end{center}
\end{titlepage}
\hypersetup{pageanchor=true}

\mysection{1}{Целое и его часть}
\idea{Вся величина --- одна единица}
В задачах на доли \textbf{целое} принимают за $1$, а в задачах на проценты --- за $100\%$. Прежде чем считать, определите, \textbf{что именно принято за целое}.
\[
1=100\%,\qquad 8\%=0{,}08,\qquad 45\%=0{,}45,\qquad
\frac35=0{,}6.
\]
\idea{Две обратные операции}
Пусть $C$ --- целое, $d$ --- его доля, $A$ --- соответствующая часть.
\[
\boxed{A=C\cdot d}\qquad
\boxed{C=\frac{A}{d}\quad(d\ne0)}.
\]
\ex{Три пятых некоторого числа равны $84$. Найдите это число.}
\[
\frac35 C=84,\qquad C=84:\frac35=84\cdot\frac53=140.
\]
Проверка: $140\cdot\frac35=84$.
\note{Ищем часть --- умножаем на её долю. Ищем целое по известной части --- делим на долю.}

\mysection{2}{Увеличение и уменьшение в процентах}
При увеличении первоначального количества на $p\%$ новая величина составляет $100\%+p\%$. При уменьшении --- $100\%-p\%$.
\[
\text{после увеличения}:\quad C\left(1+\frac p{100}\right),\qquad
\text{после уменьшения}:\quad C\left(1-\frac p{100}\right).
\]
\ex{Строители сдали $432$ тыс. м$^2$ жилья, превысив план на $8\%$. Сколько планировали сдать?}
Здесь план --- $100\%$, превышение --- $8\%$, выполненная площадь --- $108\%=1{,}08$ плана.
\begin{center}
\begin{tikzpicture}[x=0.095cm,y=1cm]
\fill[accent!12] (0,0) rectangle (100,0.67);
\fill[accent!32] (100,0) rectangle (108,0.67);
\draw[accent,line width=.8pt] (0,0) rectangle (108,0.67);
\draw[accent,line width=.7pt] (100,0)--(100,.67);
\node[font=\small] at (50,.34) {$100\%$ --- план};
\node[font=\scriptsize,anchor=south] at (104,.7) {$8\%$};
\node[font=\small,anchor=north] at (54,-.09) {$108\%$ соответствуют $432$ тыс. м$^2$};
\end{tikzpicture}
\end{center}
\[
C=432:1{,}08=400\quad\text{(тыс. м$^2$)}.
\]
Проверка: $8\%$ от $400$ --- это $32$; $400+32=432$.
\note{Фраза «на $8\%$ больше» означает $108\%$ от исходной величины, если речь идёт об увеличении исходной величины на $8\%$.}

\clearpage
\mysection{3}{Все доли приводим к одному целому}
При составлении схемы полезно подписывать сверху \textbf{доли от всей величины}, снизу --- \textbf{метры, тонны, рубли} или другие единицы измерения.

\ex{Первый участок дистанции равен $45\%$ всей длины. Второй составляет $0{,}8$ первого. Третий равен $380$ м. Найдите длину дистанции.}
Первый участок: $0{,}45$ всей дистанции. Второй участок составляет
\[
0{,}8\cdot0{,}45=0{,}36\quad\text{всей дистанции}.
\]
Тогда третий участок составляет $1-0{,}45-0{,}36=0{,}19$ всей дистанции.
\begin{center}
\begin{tikzpicture}[x=0.095cm,y=1cm]
\fill[accent!18] (0,0) rectangle (45,0.7);
\fill[accent!35] (45,0) rectangle (81,0.7);
\fill[accent!8] (81,0) rectangle (100,0.7);
\draw[accent,line width=.8pt] (0,0) rectangle (100,0.7);
\draw[accent,line width=.65pt] (45,0)--(45,.7) (81,0)--(81,.7);
\node[font=\small] at (22.5,.35) {$45\%$};
\node[font=\small] at (63,.35) {$36\%$};
\node[font=\small] at (90.5,.35) {$19\%$};
\node[font=\small,anchor=north] at (22.5,-.08) {I};
\node[font=\small,anchor=north] at (63,-.08) {II};
\node[font=\small,anchor=north] at (90.5,-.08) {III: $380$ м};
\end{tikzpicture}
\end{center}
Доля $0{,}19$ соответствует $380$ м. Поэтому
\[
C=380:0{,}19=2000\ \text{м}.
\]
Первый участок: $2000\cdot0{,}45=900$ м; второй: $900\cdot0{,}8=720$ м.
Проверка: $900+720+380=2000$ м.
\note{Доля второго участка ($0{,}8$) дана \emph{от первого}, поэтому сначала выражаем первую часть в долях целого, а затем умножаем на $0{,}8$.}

\mysection{4}{Проценты от остатка}
\ex{В первый день отгрузили $40\%$ всего зерна, во второй --- $60\%$ остатка. На третий день осталось $96$ т. Сколько зерна было первоначально?}
После первого дня осталось $1-0{,}4=0{,}6$ всего зерна. Во второй день отгрузили
\[
0{,}6\cdot0{,}6=0{,}36\quad\text{всего зерна}.
\]
На третий день пришлось
\[
1-0{,}4-0{,}36=0{,}24\quad\text{всего зерна}.
\]
Поэтому $C=96:0{,}24=400$ т.
Проверка: сначала $160$ т, затем $60\%$ от $240$ т, то есть $144$ т; осталось $400-160-144=96$ т.
\note{$60\%$ \emph{остатка} и $60\%$ \emph{первоначального количества} --- разные величины. Каждый процент должен иметь явно названное основание.}

\clearpage
\mysection{5}{Обыкновенные дроби и последовательное расходование}
С дробями действует тот же принцип: сначала выразите все необходимые части в \textbf{долях первоначального целого}. Если берут долю остатка, умножайте её на долю этого остатка.

\ex{В первый день израсходовали $\frac25$ запасов, во второй --- $\frac13$ остатка. После этого осталось $72$ кг. Сколько было сначала?}
После первого дня осталось $1-\frac25=\frac35$ запаса. Во второй день израсходовали $\frac13\cdot\frac35=\frac15$ первоначального количества. На конец второго дня осталось
\[
1-\frac25-\frac15=\frac25.
\]
Доля $\frac25$ соответствует $72$ кг, значит
\[
C=72:\frac25=72\cdot\frac52=180\ \text{кг}.
\]
Проверка: $\frac25$ от $180$ --- $72$ кг; остаток $108$ кг. Треть остатка --- $36$ кг; снова осталось $72$ кг.

\mysection{6}{Самопроверка и типичные ошибки}
\begin{itemize}
\item \textbf{Неверное целое.} Выясните, от чего считают процент: от первоначального количества, остатка или другой части.
\item \textbf{Смешение процентов и долей.} Перед вычислениями переводите проценты в десятичные дроби: $35\%=0{,}35$.
\item \textbf{Неверное действие.} Известную часть делим на её долю; долю от известного целого находим умножением.
\item \textbf{Путаница с увеличением.} После увеличения на $15\%$ получаем $115\%=1{,}15$ исходной величины.
\item \textbf{Проверка забыта.} Сумма всех частей должна совпадать с целым; повторная подстановка должна восстановить известное число.
\end{itemize}
\idea{Удобная схема данных}
В верхней строке пишите доли \emph{от одного и того же целого}, в нижней --- соответствующие числа и единицы измерения. Схему можно не рисовать, если связи между величинами очевидны.

\clearpage
\thispagestyle{plain}
\begin{center}
{\Large\color{darkaccent}\bfseries Алгоритм нахождения целого по его части}\par
\vspace{10mm}
\begin{tikzpicture}[
 >={Stealth[length=2.4mm]},line width=.75pt,
 term/.style={ellipse,draw=accent,minimum width=4cm,minimum height=8mm,align=center,fill=accent!7,font=\small},
 action/.style={rectangle,rounded corners=2pt,draw=accent,minimum width=8cm,text width=7.5cm,minimum height=12mm,align=center,font=\small,inner sep=5pt},
 decision/.style={diamond,aspect=3,draw=accent,text width=4.2cm,align=center,inner sep=2pt,font=\small},
 side/.style={rectangle,rounded corners=2pt,draw=accent,text width=3.0cm,align=center,font=\small,inner sep=6pt},
 arr/.style={->,accent}
]
\node[term] (start) {Начало};
\node[action,below=7mm of start] (whole) {Определите, что принято за целое ($1$ или $100\%$)};
\node[action,below=6mm of whole] (percent) {Переведите проценты и дроби в доли};
\node[decision,below=8mm of percent] (test) {Есть доля от остатка или другой части?};
\node[side,right=14mm of test] (rebase) {Умножьте доли, чтобы перейти к первоначальному целому};
\node[action,below=13mm of test] (same) {Запишите все части как доли первоначального целого};
\node[action,below=6mm of same] (remainder) {При необходимости найдите недостающую долю из $1$};
\node[action,below=6mm of remainder] (pair) {Найдите пару: известная часть и её доля};
\node[action,below=6mm of pair] (compute) {Целое = известная часть $:$ её доля};
\node[action,below=6mm of compute] (verify) {Проверьте ответ по условию задачи};
\node[term,below=7mm of verify] (end) {Ответ};
\draw[arr] (start)--(whole);
\draw[arr] (whole)--(percent);
\draw[arr] (percent)--(test);
\draw[arr] (test)--node[left,font=\scriptsize]{нет}(same);
\draw[arr] (test.east)--node[above,font=\scriptsize]{да}(rebase.west);
\draw[arr] (rebase.south) |- (same.east);
\draw[arr] (same)--(remainder);
\draw[arr] (remainder)--(pair);
\draw[arr] (pair)--(compute);
\draw[arr] (compute)--(verify);
\draw[arr] (verify)--(end);
\end{tikzpicture}
\end{center}

\clearpage
\mysection{7}{Тренировочный блок --- Путь А}
\textit{10 упражнений по возрастанию сложности.} Во всех сюжетных задачах определите первоначальное целое. Для задач с несколькими частями полезно сделать схему.
\idea{Базовый уровень}
\begin{enumerate}
\item Найдите $35\%$ от числа $240$.
\item Число $72$ составляет $\frac38$ неизвестного числа. Найдите его.
\item Рабочие выполнили $324$ изделия, превысив план на $8\%$. Сколько изделий было запланировано?
\item После скидки $15\%$ товар стоит $680$ руб. Сколько он стоил до скидки?
\end{enumerate}
\idea{Средний уровень}
\begin{enumerate}[resume]
\item Первый участок пути составляет $30\%$ всей дистанции, второй --- половину первого. Третий участок равен $550$ м. Найдите длину дистанции.
\item В первый день продали $25\%$ запаса яблок, во второй --- $40\%$ оставшихся. После двух дней осталось $180$ кг. Сколько яблок было изначально?
\item В первый день израсходовали $\frac38$ суммы денег, во второй --- $\frac25$ остатка. После этого осталось $150$ руб. Сколько было сначала?
\item Туристы прошли сначала $40\%$ маршрута, затем ещё $75\%$ расстояния первого перехода. Им осталось пройти $1800$ м. Найдите длину маршрута.
\end{enumerate}
\idea{Повышенный уровень}
\begin{enumerate}[resume]
\item В первый день библиотека выдала $20\%$ книг со стеллажа, во второй --- $\frac34$ оставшихся. После этого на стеллаже осталось $48$ книг. Сколько книг было сначала?
\item Первый участок беговой дистанции равен $\frac5{12}$ всей длины. Второй составляет $\frac35$ первого. Третий равен $300$ м. Найдите длину дистанции, а также первого и второго участков.
\end{enumerate}

\clearpage
\mysection{8}{Ответы к тренировочному блоку}
\renewcommand{\arraystretch}{1.45}
\begin{center}
\begin{tabularx}{.91\linewidth}{@{}cX@{}}
\toprule
№ & Ответ \\
\midrule
1 & $84$ \\
2 & $192$ \\
3 & $300$ изделий \\
4 & $800$ руб. \\
5 & $1000$ м \\
6 & $400$ кг \\
7 & $400$ руб. \\
8 & $6000$ м \\
9 & $240$ книг \\
10 & Вся дистанция --- $900$ м; первый участок --- $375$ м; второй --- $225$ м \\
\bottomrule
\end{tabularx}
\end{center}
\vspace{9mm}
\idea{Итоговая памятка}
\begin{tabularx}{\linewidth}{@{}>{\raggedright\arraybackslash}p{.35\linewidth}X@{}}
\toprule
Сигнал в задаче & Полезное действие\\
\midrule
«Составляет $p\%$» & Переведите процент в долю $p/100$.\\
«На $p\%$ больше» & Используйте долю $1+p/100$ от исходной величины.\\
«На $p\%$ меньше» & Используйте долю $1-p/100$ от исходной величины.\\
«Процент остатка» & Сначала найдите долю остатка, затем перемножьте доли.\\
«Известна часть» & Найдите её долю от первоначального целого и разделите часть на эту долю.\\
\bottomrule
\end{tabularx}
\vfill
{\small\color{muted}Лёвин Артём Александрович --- эксклюзивно для Нади Клименко}
\end{document}
